% The actual content of the blueprint lives here. It is shared by both the
% web version (web.tex) and the print version (print.tex).

\section{The red/black card game}

We study the optimal-stopping equity $e(r,b)$ of the game with $r$ red and $b$
black cards, where a red draw pays $+1$, a black draw pays $-1$, and the player
may stop at any time.

\begin{definition}[Equity]\label{def:e}
\lean{e}\leanok
Define $e:\mathbb N\times\mathbb N\to\mathbb Q$ by
\[
  e(0,b)=0,\qquad e(r,0)=r,
\]
\[
  e(r+1,b+1)=\max\!\Big(0,\ \tfrac{r+1}{r+b+2}\big(e(r,b+1)+1\big)
                          +\tfrac{b+1}{r+b+2}\big(e(r+1,b)-1\big)\Big).
\]
\end{definition}

We record three facts already established in Lean.

\begin{lemma}[Nonnegativity]\label{lem:zero_le_e}
\lean{zero_le_e} \leanok
\uses{def:e}
$e(r,b)\ge 0$ for all $r,b$.
\end{lemma}
\begin{proof}\leanok
Formalised in Lean (\texttt{zero\_le\_e}).
\end{proof}

\begin{lemma}[Lower bound]\label{lem:sub_le_e}
\lean{sub_le_e}\leanok
\uses{def:e}
$r-b\le e(r,b)$ for all $r,b$, the inequality being read in $\mathbb Q$ after
coercion.
\end{lemma}
\begin{proof}\leanok
Formalised in Lean (\texttt{sub\_le\_e}).
\end{proof}

\begin{lemma}[Positivity on the diagonal]\label{lem:e_pos}
\lean{e_pos_of_pos}\leanok
\uses{def:e}
$e(n,n)>0$ for $n>0$.
\end{lemma}
\begin{proof}\leanok
Formalised in Lean (\texttt{e\_pos\_of\_pos}).
\end{proof}

\section{Goal}

\begin{theorem}[\texttt{question}]\label{thm:question}
\lean{question}\leanok
\uses{thm:super_6_4, prop:comparison, lem:zero_le_e}
If $b>r+4\sqrt r$, where $\sqrt r$ is the real square root and the inequality is
read in $\mathbb R$ after casting $r,b$, then $e(r,b)=0$ --- for \emph{all} $r,b$,
with no lower bound on $r$.
\end{theorem}

The constant $4$ is not special: it is the width constant of one particular barrier
we choose to formalise (Section~\ref{sec:instance}). The proof is an upper-bound
(supersolution) argument valid for any pair $(K,c)$ meeting the constraints of
Theorem~\ref{thm:general}; a \emph{smaller} $c$ gives a \emph{tighter} zero region,
and $c$ can be pushed down to roughly $1+\sqrt2$ before the constraints become
infeasible. (An earlier version of this file used the deliberately generous
placeholder $c=137$; nothing depended on that value beyond it being large enough.)

\section{Recursion and the reduction to an upper bound}

\begin{lemma}[Unfolded recursion]\label{lem:recursion}
\lean{e_succ_succ}\leanok
\uses{def:e}
For all $r,b\ge 1$, writing $N=r+b$,
\[
  e(r,b)=\max\!\Big(0,\ \tfrac{r}{N}\big(e(r-1,b)+1\big)
                       +\tfrac{b}{N}\big(e(r,b-1)-1\big)\Big).
\]
\end{lemma}
\begin{proof}\leanok
Write $r=r'+1$, $b=b'+1$ and apply Definition~\ref{def:e}; then $N=r'+b'+2=r+b$,
$e(r',b'+1)=e(r-1,b)$ and $e(r'+1,b')=e(r,b-1)$.
\end{proof}

\begin{lemma}[Reduction]\label{lem:reduction}
\lean{e_eq_zero_of_inner_nonpos}\leanok
\uses{lem:recursion, lem:zero_le_e}
For $r,b\ge 1$, if
\[
  \tfrac{r}{N}\big(e(r-1,b)+1\big)+\tfrac{b}{N}\big(e(r,b-1)-1\big)\le 0,
\]
then $e(r,b)=0$.
\end{lemma}
\begin{proof}\leanok
By Lemma~\ref{lem:recursion} the right-hand side of the $\max$ is $\le 0$, so the
$\max$ equals $0$; combined with $e(r,b)\ge 0$ this gives $e(r,b)=0$.
\end{proof}

\section{The comparison (supersolution) principle}

\begin{definition}[Supersolution]\label{def:super}
\lean{IsSupersolution}\leanok
\uses{def:e}
A real-valued function $\varphi:\mathbb N\times\mathbb N\to\mathbb R$ is a
\emph{supersolution} if
\begin{enumerate}
  \item $\varphi(r,b)\ge 0$ for all $r,b$;
  \item $\varphi(0,b)=0$ and $\varphi(r,0)\ge r$;
  \item for all $r,b\ge1$, writing $N=r+b$,
  \[
    \tfrac{r}{N}\big(\varphi(r-1,b)+1\big)+\tfrac{b}{N}\big(\varphi(r,b-1)-1\big)
      \ \le\ \varphi(r,b). \tag{SS}
  \]
\end{enumerate}
\end{definition}

\begin{proposition}[Comparison]\label{prop:comparison}
\lean{e_le_of_supersolution}\leanok
\uses{lem:recursion, def:super}
If $\varphi$ is a supersolution, then $e(r,b)\le\varphi(r,b)$ for all $r,b$ (the
rational $e(r,b)$ being cast into $\mathbb R$).
\end{proposition}
\begin{proof}\leanok
Strong induction on $r+b$.
If $r=0$ then $e(0,b)=0\le\varphi(0,b)$ since $\varphi\ge0$.
If $b=0$ then $e(r,0)=r\le\varphi(r,0)$ by (2).
For $r,b\ge1$, Lemma~\ref{lem:recursion} gives
$e(r,b)=\max(0,\ \tfrac{r}{N}(e(r-1,b)+1)+\tfrac{b}{N}(e(r,b-1)-1))$;
the coefficients $r/N,\ b/N$ are nonnegative, so by the induction hypothesis and
(SS) the inner expression is $\le\varphi(r,b)$, and since $0\le\varphi(r,b)$ the
maximum is $\le\varphi(r,b)$.
\end{proof}

Thus, to prove $e(r,b)=0$ on a region it suffices to exhibit a supersolution
$\varphi$ that \emph{vanishes} there: then $0\le e(r,b)\le\varphi(r,b)=0$.

\section{Lower bounds via subsolutions}

Dual to the comparison principle is a lower-bound principle. A
\emph{subsolution} reverses the supersolution inequalities and drops the
nonnegativity requirement; where a supersolution dominates $e$ from above, a
subsolution is dominated by $e$ from below.

\begin{definition}[Subsolution]\label{def:subsolution}
\lean{IsSubsolution}\leanok
\uses{def:e}
A real-valued function $\psi:\mathbb N\times\mathbb N\to\mathbb R$ is a
\emph{subsolution} if
\begin{enumerate}
  \item $\psi(0,b)\le 0$ for all $b$;
  \item $\psi(r,0)\le r$ for all $r$;
  \item for all $r,b\ge1$, writing $N=r+b$,
  \[
    \psi(r,b)\ \le\
    \tfrac{r}{N}\big(\psi(r-1,b)+1\big)+\tfrac{b}{N}\big(\psi(r,b-1)-1\big).
    \tag{SS$'$}
  \]
\end{enumerate}
Because $e=\max(0,\cdot)$, no truncation is needed in (SS$'$): it suffices that
$\psi$ lie below its own ``continue value''.
\end{definition}

\begin{proposition}[Lower comparison]\label{prop:comparison_lower}
\lean{e_ge_of_subsolution}\leanok
\uses{lem:recursion, def:subsolution}
If $\psi$ is a subsolution, then $\psi(r,b)\le e(r,b)$ for all $r,b$ (the rational
$e(r,b)$ being cast into $\mathbb R$).
\end{proposition}
\begin{proof}\leanok
Strong induction on $r+b$, mirroring Proposition~\ref{prop:comparison}.
If $r=0$ then $\psi(0,b)\le 0=e(0,b)$ by (1). If $b=0$ then
$\psi(r,0)\le r=e(r,0)$ by (2). For $r,b\ge1$, the coefficients $r/N,\ b/N$ are
nonnegative, so by the induction hypothesis the continue value of $\psi$ is at most
that of $e$; with (SS$'$) and $e(r,b)=\max(0,\cdot)\ge(\cdot)$ this gives
$\psi(r,b)\le e(r,b)$.
\end{proof}

\begin{lemma}[Linear subsolution]\label{lem:sub_subsolution}
\lean{sub_isSubsolution}\leanok
\uses{def:subsolution}
The function $\psi(r,b)=r-b$ is a subsolution: it satisfies (SS$'$) with equality
(it is the value of the ``draw everything'' strategy).
\end{lemma}
\begin{proof}\leanok
Conditions (1),(2) are immediate. For (SS$'$), the equality
$r\big((r-b)-(r-1-b)\big)+b\big((r-b)-(r-b+1)\big)=r-b$ makes the continue value of
$\psi$ equal to $\psi(r,b)$.
\end{proof}

Chaining Lemma~\ref{lem:sub_subsolution} through
Proposition~\ref{prop:comparison_lower} recovers the lower bound
$r-b\le e(r,b)$ of Lemma~\ref{lem:sub_le_e}, now as an instance of the
subsolution principle. Combined with the upper bound of
Theorem~\ref{thm:question}, this gives a two-sided picture: $e$ vanishes for
$b>r+4\sqrt r$ and is bounded below by $r-b$ throughout.

\section{The barrier and the general supersolution theorem}

\begin{definition}[Quadratic barrier $\varphi_{K,c}$]\label{def:phiBar}
\lean{phiBar}\leanok
\uses{def:e}
For real parameters $K,c$ (a normalisation and a width constant) set, with
$D=r+c\sqrt r-b$,
\[
  \varphi_{K,c}(r,b)=
  \begin{cases}
    0, & D\le 0 \quad(\text{far field}),\\[2pt]
    D^2/(K\sqrt r), & D>0,\ r\le b \quad(\text{layer}),\\[2pt]
    (r-b)+\tfrac{c^2}{K}\sqrt r, & D>0,\ b<r \quad(\text{bulk}).
  \end{cases}
\]
\end{definition}

\begin{theorem}[General supersolution theorem]\label{thm:general}
\lean{phiBar_isSupersolution_of}\leanok
\uses{def:super, def:phiBar, lem:layerform, prop:ss_far, prop:ss_bulk, prop:ss_hard}
Suppose $(K,c)$ satisfy the four polynomial side-conditions
\[
  1\le K,\qquad K\le 2c,\qquad 1\le c(K-1),\qquad
  K^2-8Kc+2K+4c^2+16c+13\ \le\ 0,
\]
together with the \emph{layer inequality}: for all $s\ge1$, $w\ge0$, $t\ge0$ with
$t^2=s^2+1$,
\[
  0\ \le\ s(ct-w)^2(2s^2+w+2)-(s^2+1)\big(t(cs-w-1)^2+Kst\big)
          -(s^2+w+1)\big(s(ct-w+1)^2-Kst\big). \tag{L}
\]
Then $\varphi_{K,c}$ is a supersolution.
\end{theorem}

The last polynomial condition already forces $c\ge 1+\sqrt2$; it is exactly the
requirement that the $r=0$ base case $3b^2+(K-4c+1)b+(c^2-2c-1)\ge0$ hold for all
$b$ (nonnegative discriminant of the admissible-$K$ window). The remaining three
are the positivity of $\varphi_{K,c}$ ($K>0$), the bulk~$\le$~layer-form bound
($K\le2c$), and the far-field bound ($1\le c(K-1)$). The genuinely
$(K,c)$-dependent content is the layer inequality~(L).

\paragraph{Why a $\sqrt r$ layer.}
Two observations constrain the shape of $\varphi$. First, $r-b$ satisfies (SS) with
equality: $r((r-b)-(r-1-b))+b((r-b)-(r-b+1))=r-b$; this is the value of the naive
``draw everything'' strategy. Second, $\max(0,r-b)$ is \emph{not} a supersolution:
at the diagonal $b=r$ one has $\varphi(r,r)=0$ but $\varphi(r,r-1)=1$, and (SS)
there fails. The defect is governed by the known growth $e(r,r)=\Theta(\sqrt r)$,
so a dominating supersolution must satisfy $\varphi(r,r)=\Omega(\sqrt r)$ while
vanishing at $b=r+c\sqrt r$: a diffusive boundary layer of width $\sqrt r$, the
origin of the $\sqrt r$ term in the threshold.

\paragraph{The continuum window (heuristic).}
Expanding the two discrete differences of the layer $\varphi_{K,c}=D^2/(K\sqrt r)$,
with $s=\sqrt r$ and $x=D/s\in[0,c]$,
\[
  r\,\Delta_r\varphi+b\,\Delta_b\varphi-(r-b)\;=\;s\,P(x)+O(1)+O(1/s),
  \qquad
  P(x)=\tfrac{3}{2K}x^2-\big(\tfrac cK+1\big)x+\big(c-\tfrac2K\big).
\]
So to leading order the layer is a supersolution iff $P\ge0$ on $[0,c]$, whose
interior minimum at $x_*=(c+K)/3$ gives
\[
  P(x_*)\ge0\iff K^2-4cK+c^2+12\le0
  \iff K\in\big[\,2c-\sqrt{3c^2-12},\ 2c+\sqrt{3c^2-12}\,\big],
\]
nonempty exactly when $c\ge2$. The rigorous discrete constraints of
Theorem~\ref{thm:general} are a slightly stronger (finite-$r$) version of this
window, pushing the feasibility floor from $c\ge2$ up to $c\ge1+\sqrt2$.

\paragraph{The continuum profile (heuristic).}
To second order the (SS) inequality is, in the continuum, an obstacle problem with
obstacle $0$ whose continuation equation
$r\varphi_r+b\varphi_b-\tfrac12(r\varphi_{rr}+b\varphi_{bb})=r-b$, under the
self-similar ansatz $\varphi=\sqrt r\,G(\xi)$ with $\xi=(b-r)/\sqrt r$, reduces near
the diagonal to the Hermite-type ODE $G''-\xi G'-\tfrac12G=\xi$ with smooth-fit
data $G(c)=G'(c)=0$ and bulk matching $G\sim-\xi$. The genuine boundary constant is
a fixed $O(1)$ number; the quadratic layer is an explicit elementary
over-estimate for which the \emph{discrete} (SS) can be checked.

\paragraph{Region lemmas.}
The supersolution inequality (SS) for $\varphi_{K,c}$ is proved by splitting on
which piece each of the three involved points lands in.

\begin{proposition}[Far field]\label{prop:ss_far}
\lean{phiBar_ss_far_of}\leanok
\uses{def:phiBar}
If $1\le K$, $1\le c$, $1\le c(K-1)$ and the target point is in the far field, then
(SS) holds; the left-hand side is in fact $\le0$.
\end{proposition}
\begin{proof}\leanok Formalised: after rewriting the vanishing points, the one
surviving layer term is controlled by $D\le1$ and $c\sqrt{r+1}\le b-r$, closed by a
single \texttt{nlinarith} over the cleared formulas. \end{proof}

\begin{proposition}[Bulk]\label{prop:ss_bulk}
\lean{phiBar_ss_bulk_of}\leanok
\uses{def:phiBar}
If $0<K$ and $b+1<r$, then (SS) holds; it collapses to $\sqrt r\le\sqrt{r+1}$, the
$c^2/K$ and $N$ factors cancelling.
\end{proposition}
\begin{proof}\leanok Formalised (\texttt{phiBar\_ss\_bulk\_of}). \end{proof}

\begin{proposition}[Layer-form bound]\label{lem:layerform}
\lean{phiBar_le_layerform_of}\leanok
\uses{def:phiBar}
If $0<K\le 2c$ and $R\ge1$, then $\varphi_{K,c}(R,b)\le(R+c\sqrt R-b)^2/(K\sqrt R)$;
in the bulk the cleared difference is $(R-b)\big((R-b)+(2c-K)\sqrt R\big)\ge0$.
\end{proposition}
\begin{proof}\leanok Formalised (\texttt{phiBar\_le\_layerform\_of}). \end{proof}

\begin{proposition}[Hard (layer/seam) case]\label{prop:ss_hard}
\lean{phiBar_ss_hard_of}\leanok
\uses{def:phiBar, lem:layerform}
Assuming the $r=0$ discriminant condition and the layer inequality (L) of
Theorem~\ref{thm:general}, (SS) holds on the layer $r\le b$ and the diagonal seam
$b=r-1$.
\end{proposition}
\begin{proof}\leanok
Three sub-cases: the $r=0$ base reduces to $3b^2+(K-4c+1)b+(c^2-2c-1)\ge0$, which
follows from the discriminant condition; the seam $b=r-1$ collapses to
$c^2(b+2)(\sqrt{b+2}-\sqrt{b+1})\ge0$; and the main layer $b\ge r$ is
$\varphi\le$~layer-form (Proposition~\ref{lem:layerform}) followed by the divided
layer inequality (L).
\end{proof}

\section{The concrete instance $(K,c)=(6,4)$ and its certificate}
\label{sec:instance}

We instantiate Theorem~\ref{thm:general} with $K=6$, $c=4$. The four side-conditions
hold ($6\le8$, $1\le4\cdot5=20$, and $6^2-8\cdot6\cdot4+2\cdot6+4\cdot4^2+16\cdot4+13=-3\le0$),
so it remains to verify the layer inequality (L). We clear the square root: writing
$s=\sqrt r$, $t=\sqrt{r+1}$ ($t^2=s^2+1$), $w=b-r$, (L) has the shape $A-tB_{+}$ with
$A,B_{+}\ge0$, hence is equivalent to the square-root-free $A^2\ge(s^2+1)B_{+}^2$,
i.e. $G(s,w)\ge0$.

\begin{lemma}[Square-root-free core]\label{lem:G}
\lean{hard_G_nonneg_6_4}\leanok
For $s\ge1$ and all real $w$, $G(s,w)\ge0$, via the closed-form Positivstellensatz
identity
\[
  4\,c_4\,H_2\,G \;=\; H_2\,(2c_4w^2+c_3w)^2 \;+\; (H_2w+2c_4c_1)^2 \;+\; 4\,c_4\,H_3,
\]
with $c_4=3s^4+6s^2-1>0$, $H_2>0$, $H_3\ge0$.
\end{lemma}
\begin{proof}\leanok
The displayed identity is an exact \texttt{ring} identity; nonnegativity of the
right-hand side follows from the sign facts $c_4>0$, $H_2>0$, $H_3\ge0$ (each a
univariate polynomial inequality in $s$ whose coefficients are all nonnegative
after the shift $s=1+u$, $u\ge0$, discharged by \texttt{nlinarith}), and dividing
by $4c_4H_2>0$ gives $G\ge0$.
\end{proof}

\begin{lemma}[Layer inequality at $(6,4)$]\label{lem:layercore}
\lean{layer_core_6_4}\leanok
\uses{lem:G}
The layer inequality (L) holds for $K=6$, $c=4$.
\end{lemma}
\begin{proof}\leanok
Combine $G\ge0$ (Lemma~\ref{lem:G}) with the sign facts $A,B_{+}\ge0$, $t\ge0$ to
remove the square root, then rewrite $A-tB_{+}$ into the (L) form via $t^2=s^2+1$.
\end{proof}

\begin{theorem}[$\varphi_{6,4}$ is a supersolution]\label{thm:super_6_4}
\lean{phiBar_isSupersolution_6_4}\leanok
\uses{thm:general, lem:layercore}
$\varphi_{6,4}$ is a supersolution.
\end{theorem}
\begin{proof}\leanok
Apply Theorem~\ref{thm:general} with the four side-conditions checked numerically
and the layer inequality supplied by Lemma~\ref{lem:layercore}.
\end{proof}

Chaining Theorem~\ref{thm:super_6_4} $\to$ Proposition~\ref{prop:comparison} and
evaluating $\varphi_{6,4}$ in the far field gives Theorem~\ref{thm:question}: for
$b>r+4\sqrt r$ we have $0\le e(r,b)\le\varphi_{6,4}(r,b)=0$. The whole chain compiles
with no \texttt{sorry}.

\section{Below-threshold positivity (complementary lower bound)}

We now show that below a matching threshold $b<r+k\sqrt r$ the equity is
\emph{strictly positive}. The two-dimensional positivity is reduced to a purely
one-dimensional statement (a lower bound on the diagonal) by a Lipschitz argument;
an elementary subsolution barrier is \emph{not} available for this direction
(a subharmonic correction cannot also be a positive bump vanishing at both ends).

\begin{lemma}[Bounded differences]\label{lem:lipschitz}
\lean{e_lipschitz}\leanok
\uses{lem:recursion, lem:sub_le_e, lem:zero_le_e}
For all $r,b$, one has $e(r,b)\le e(r,b+1)+1$ and $e(r+1,b)\le e(r,b)+1$: one extra
black card lowers the equity by at most $1$, and one extra red card raises it by at
most $1$.
\end{lemma}
\begin{proof}\leanok
Simultaneous strong induction on $r+b$. At $(r,b)$ the two claims each reduce, via
the recursion (Lemma~\ref{lem:recursion}) and clearing the denominator $N$, to the
two claims at $(r-1,b)$ / $(r,b-1)$ of the previous level; the base cases use
$e(r,0)=r$ and $r-1\le e(r,1)$ (Lemma~\ref{lem:sub_le_e}).
\end{proof}

\begin{lemma}[Diagonal iteration]\label{lem:diag_iter}
\lean{e_diag_lower}\leanok
\uses{lem:lipschitz}
For all $r$ and $j$, $e(r,r+j)\ge e(r,r)-j$.
\end{lemma}
\begin{proof}\leanok
Iterate the black-Lipschitz bound $e(r,b+1)\ge e(r,b)-1$ (Lemma~\ref{lem:lipschitz})
$j$ times starting from the diagonal.
\end{proof}

\begin{theorem}[Conditional below-threshold positivity]\label{thm:cond_positivity}
\lean{e_pos_below_threshold_of_diag}\leanok
\uses{lem:diag_iter, lem:sub_le_e}
Fix $k>0$ and suppose the diagonal lower bound $e(r,r)\ge k\sqrt r$ holds. Then
$e(r,b)>0$ for every $b$ with $b<r+k\sqrt r$.
\end{theorem}
\begin{proof}\leanok
If $b<r$ then $e(r,b)\ge r-b>0$ by Lemma~\ref{lem:sub_le_e}. If $b\ge r$, write
$b=r+j$; Lemma~\ref{lem:diag_iter} and the hypothesis give
$e(r,b)\ge e(r,r)-j\ge k\sqrt r-(b-r)>0$.
\end{proof}

It therefore remains to prove the diagonal lower bound $e(r,r)\ge k\sqrt r$. This is
the genuine $\sqrt r$-diffusion fact. Rather than the ballot/reflection route (which
Mathlib does not support), we obtain it by a \emph{reflection-free} argument: we
write down the value of one explicit strategy in closed form and verify it against
the recursion by binomial algebra.

\section{The diagonal $\sqrt r$ lower bound (reflection-free)}

Consider the strategy ``keep drawing until the number of reds drawn exceeds the
number of blacks drawn by $L$, otherwise draw to the end''. Started from the
diagonal the surplus equals $b-r$, so this strategy is a function of the state
$(r,b)$, and its value has the closed form
\[
  W_L(r,b) = (r-b) + L\cdot P_L(r,b),\qquad
  P_L(r,b) = \binom{r+b}{\,b-L}\big/\binom{r+b}{\,b},
\]
where $P_L(r,b)$ is the probability the surplus ever reaches $L$ (and $P_L=0$ when
$b<L$).

\begin{definition}[Reach probability]\label{def:PL}
\lean{PL}\leanok
$P_L(r,b)=\binom{r+b}{b-L}/\binom{r+b}{b}$ for $b\ge L$, and $0$ otherwise.
\end{definition}

\begin{lemma}[$P_L$ is a martingale]\label{lem:PL_mart}
\lean{PL_mart}\leanok
\uses{def:PL}
For $r,b\ge1$, $P_L(r,b)=\tfrac{r}{N}P_L(r-1,b)+\tfrac{b}{N}P_L(r,b-1)$, $N=r+b$.
\end{lemma}
\begin{proof}\leanok
Pure binomial algebra: Pascal's rule
$\binom{r+b}{b-L}=\binom{r+b-1}{b-L}+\binom{r+b-1}{b-1-L}$ together with the
multiplicative identities $r\binom{r+b}{b}=(r+b)\binom{r+b-1}{b}$ and
$b\binom{r+b}{b}=(r+b)\binom{r+b-1}{b-1}$. No reflection principle is used.
\end{proof}

\begin{lemma}[Strategy value satisfies the recursion]\label{lem:wL_mart}
\lean{wL_mart}\leanok
\uses{def:PL, lem:PL_mart}
$W_L(r,b)=\tfrac{r}{N}\big(W_L(r-1,b)+1\big)+\tfrac{b}{N}\big(W_L(r,b-1)-1\big)$ for
$r,b\ge1$; and $W_L(r,r+L)=0$.
\end{lemma}
\begin{proof}\leanok
The $(r-b)$ part satisfies the recursion with equality; the $L\,P_L$ part is
Lemma~\ref{lem:PL_mart}. The boundary value uses $\binom{r+b}{b-L}=\binom{r+b}{b}$
at $b=r+L$ (choose-symmetry).
\end{proof}

\begin{proposition}[Domination]\label{prop:wL_le_e}
\lean{wL_le_e}\leanok
\uses{lem:wL_mart, lem:zero_le_e, lem:sub_le_e}
$W_L(r,b)\le e(r,b)$ for all $r,b$ with $b\le r+L$.
\end{proposition}
\begin{proof}\leanok
Strong induction on $r+b$. At surplus exactly $L$, $W_L=0\le e$; in the continue
region $W_L$ equals its continue value (Lemma~\ref{lem:wL_mart}), which by the
induction hypothesis and monotone coefficients is $\le$ that of $e$, hence
$\le\max(0,\cdot)=e$.
\end{proof}

\begin{lemma}[Diagonal value]\label{lem:e_diag_ge}
\lean{e_diag_ge}\leanok
\uses{prop:wL_le_e}
For $L\le n$, $e(n,n)\ge L\binom{2n}{n-L}/\binom{2n}{n}$.
\end{lemma}
\begin{proof}\leanok Take $b=n\le n+L$ in Proposition~\ref{prop:wL_le_e}. \end{proof}

\begin{lemma}[Binomial ratio bound]\label{lem:ratio_lb}
\lean{ratio_lb}\leanok
For $L\le n$, $\binom{2n}{n-L}/\binom{2n}{n}\ge 1-L^2/n$.
\end{lemma}
\begin{proof}\leanok
Induction on $L$; the step reduces to $2L^3+3L^2+3L+1\ge0$.
\end{proof}

\begin{theorem}[Diagonal $\sqrt r$ bound]\label{thm:e_diag_sqrt}
\lean{e_diag_sqrt}\leanok
\uses{lem:e_diag_ge, lem:ratio_lb}
$e(r,r)\ge\tfrac18\sqrt r$ for all $r$.
\end{theorem}
\begin{proof}\leanok
Combining Lemmas~\ref{lem:e_diag_ge} and~\ref{lem:ratio_lb} gives
$e(n,n)\ge L(1-L^2/n)$; choosing $L=\lfloor\sqrt n\rfloor/2$ yields
$e(n,n)^2\ge n/64$ for $n\ge4$ (small cases by direct bound), i.e.
$e(n,n)\ge\sqrt n/8$.
\end{proof}

\begin{theorem}[Below-threshold positivity]\label{thm:e_pos_below}
\lean{e_pos_below}\leanok
\uses{thm:cond_positivity, thm:e_diag_sqrt}
$e(r,b)>0$ whenever $b<r+\tfrac18\sqrt r$.
\end{theorem}
\begin{proof}\leanok
Theorem~\ref{thm:e_diag_sqrt} supplies the hypothesis of
Theorem~\ref{thm:cond_positivity} with $k=\tfrac18$.
\end{proof}

Together with Theorem~\ref{thm:question} this is a two-sided result: $e(r,b)=0$ for
$b>r+4\sqrt r$ and $e(r,b)>0$ for $b<r+\tfrac18\sqrt r$.
